\section{BACKGROUND AND RELATED WORK}\label{sec:related}

This work draws on four literatures: LLM agents for cloud operations, the evaluation validity of LLM
agents, runtime governance of agent actions, and the non-LLM automation those agents must be measured
against. We cover each in turn and close by stating how the present study differs.

\subsection{The Base Architecture}
Kirubakaran et~al.~\cite{kirubakaran2025governing} describe a data plane, a policy plane, and an
agentic control plane. Four bounded agents share the control plane and produce only a validated
action proposal. Their evaluation compares the system against a static-orchestration baseline across
injected failures. Our implementation follows that same decomposition; wherever the description
leaves a choice unspecified, we record our decision and its alternatives in the accompanying
artifact.

\subsection{Heuristic Automation as the Baseline}
Workflow orchestrators such as Airflow, Prefect, and Dagster are execution-centric: they schedule,
retry, and alert, without reasoning about a failure their author never anticipated. Automation
driven only by resource signals already accomplishes a great deal at production scale. Google's
Autopilot~\cite{rzadca2020autopilot}, a heuristic autoscaler covering a large share of Google's
fleet, uses no LLM at all. This kind of automation is a legitimate comparator, not a strawman. An
empirical study of real data pipelines catalogs 41 failure factors~\cite{foidl2024datapipeline}, and
comparing against static orchestration alone cannot show whether an agent system handles them better
than cheap automation does. Two idealized heuristic baselines in our own evaluation
(Section~\ref{sec:results}) beat the full agent system on recovery time and manual interventions.

\subsection{LLM Agents for Cloud Operations}
Reasoning-and-acting loops~\cite{yao2022react} are the general pattern behind agents that observe a
system and propose actions. In root-cause analysis, RCACopilot~\cite{chen2024rcacopilot} and
RCAgent~\cite{wang2023rcagent} apply tool-augmented LLMs to production incidents at Microsoft and
Alibaba Cloud. Those systems stop at a recommended diagnosis, while the design studied here also
\emph{executes} an action, which is exactly what makes a policy gate necessary.
OpsAgent~\cite{luo2025opsagent}, a multi-agent incident-management system running in production at
Lenovo, shows industrial practice heading the same way.

\subsection{Evaluation Validity for LLM Agents}
Our central finding, a mock and a live evaluation of the identical system diverging by orders of
magnitude (Section~\ref{sec:livemock}), fits inside a growing literature on the sim-to-real gap for
LLM agents. The closest theoretical framing decomposes the
gap into observation, action, transition, and reward mismatches within a unified MDP view of
foundation-model agents. The empirical evidence agrees: one study finds
that LLM-simulated environments are ``excessively cooperative'' and inflate agent success relative to
grounded real-user data. Another study~\cite{gupta2026reliabilitybench} measures a reliability drop once the
same agents move from clean benchmark conditions to production-like stress such as task
perturbation, tool faults, and rate limiting. TheAgentCompany~\cite{xu2024agentcompany}, a
high-fidelity simulated workplace, finds that even its best agent completes only about 30\% of
consequential tasks, which marks benchmark realism as a problem distinct from the mock-versus-live
axis. Benchmark validity is questioned more broadly too: a survey of over 1{,}300 agentic benchmark papers~\cite{rystrom2026agentbenchmarks}
finds that none meet basic procedural-validity criteria for public-sector deployment. We contribute a
controlled mock-versus-live comparison of one system through one code path.

\subsection{Graduated Autonomy and Runtime Policy}
The trust core described in Section~\ref{sec:system}, with its shadow, approval, and autonomous
modes, a durable kill switch, and a blast-radius cap, traces back to the game-theoretic treatment of
shutdown incentives~\cite{hadfieldmenell2016offswitch}. A related framework~\cite{zheng2026separating} casts
the same idea in concrete governance terms by separating an agent's \emph{allowed} autonomy level
from its \emph{capability} level, and demonstrates the split on an enterprise data-engineering
agent whose structure closely resembles the problem studied here.

The policy gate itself runs on Open Policy Agent~\cite{opa} and its Rego language. The base design's
safety case is that policy, not the model, bounds what can happen, and that case is only as strong as
the policy's behavior on inputs its authors never considered. One critique~\cite{joshi2026deontic}
argues that Rego-style languages are not expressive enough for agentic governance, a claim our
adversarial corpus (Section~\ref{sec:results}) puts to an empirical test. ActGov~\cite{zhang2026actgov}
validates every proposed tool call against an SMT-verified policy engine before execution, which is
functionally the same per-action gate that our \texttt{ProposedAction} contract and OPA policies
implement.

\subsection{Adversarial Evaluation}
Asserting that a policy bound holds is not the same as testing it. AgentDojo
\cite{debenedetti2024agentdojo} and the Agent Security Bench~\cite{zhang2024asb}, which reports
attack success rates up to 84.3\% against weak defenses, come closest among extensible benchmarks to
the corpus used here. The difference lies in the oracle: ours is an independently written
specification of the policy rather than the policy engine under test, so a containment figure does
not depend on the policy grading itself.

\subsection{Decision Quality and Cross-Model Generalization}
Resolving an incident quickly is not the same as resolving it correctly. AIOpsLab
\cite{chen2025aiopslab} evaluates cloud agents across detection, localization, and mitigation rather
than on latency alone, and OpenSec~\cite{barnes2026opensec} finds that frontier models detect
incidents well but calibrate poorly, with false-positive containment-action rates of 82.5\% and 45\%
for two current models. Our correct-mitigation metric (Section~\ref{sec:method}) rests on that same
distinction.

The base paper states that the GPT, Claude, and Gemini model families worked as interchangeable
backends, yet it reports no cross-model results. Evidence from elsewhere counsels caution. One study~\cite{sanna2025crossllm} measures a 43.4-point
accuracy drop when a behavioral backdoor detector moves from in-distribution to cross-LLM evaluation
across six production models, and another~\cite{zhang2026towardagentic} shows that on identical input, agent behavior shifts
independently with model version, open- versus closed-weight status, serving provider, and system
prompt, with weight-release status a better predictor of behavior than model identity. The
cross-model harness described in Section~\ref{sec:system} exists precisely to test that claim, but
its live sweep was not run in this campaign (Section~\ref{sec:discussion}).

\subsection{Fault Injection and Statistics}
Deliberately injecting faults to test resilience is established practice~\cite{basiri2016chaos}. Our
seeded injector reproduces four scenarios exactly, namely schema drift, upstream delay, resource
contention, and ingress burst. For inference we rely on bootstrap resampling~\cite{efron1979}, the
Mann--Whitney rank-sum~\cite{mann1947} and Wilcoxon signed-rank~\cite{wilcoxon1945} tests, Holm's
step-down correction~\cite{holm1979}, Cliff's $\delta$ as an ordinal effect size~\cite{cliff1993},
and Wilson score intervals~\cite{wilson1927} for proportions near one, where the normal approximation
collapses to a zero-width interval.

\subsection{Positioning}
The base paper positions policy-bounded agentic control as its advantage over unconstrained
LLM-driven automation, yet it asserts that the bound holds without testing it against an adversary,
offers no cross-model results for its model-agnostic claim, and compares only against static
orchestration. We test the first claim adversarially, benchmark against explicit, idealized non-LLM
baselines, and report plainly that the harness for the second exists but has not been run live.
Prior work on the sim-to-real gap already establishes that agents can fail under live or adversarial
conditions. What this study adds is a demonstration, inside one controlled campaign, that the same
system run through the same code path reaches qualitatively different conclusions depending on
whether the model in the loop is mocked or live.
